% 可独立编译的示例。使用 Tectonic 或 XeLaTeX。
\documentclass[11pt,a4paper]{ctexart}

\usepackage[margin=2.5cm]{geometry}
\usepackage{amsmath,amssymb,amsthm}

\newcommand{\studentname}{示例同学}
\newcommand{\studentid}{2026000000}
\newcommand{\coursename}{中级微观经济学}
\newcommand{\assignmentnumber}{1}
\newcommand{\assignmentdate}{2026-09-28}

\newcommand{\R}{\mathbb{R}}
\newcommand{\E}{\mathbb{E}}
\newcommand{\Var}{\operatorname{Var}}
\newcommand{\Cov}{\operatorname{Cov}}
\newcommand{\dd}{\mathop{}\!\mathrm{d}}

\newtheorem{problem}{题目}
\renewcommand{\proofname}{证明}
\newenvironment{solution}
  {\par\noindent\textbf{解答：}\ }
  {\par\medskip}

\begin{document}
\begin{center}
  {\Large\bfseries \coursename}\par\medskip
  作业 \assignmentnumber\par\medskip
  姓名：\studentname\quad 学号：\studentid\par
  日期：\assignmentdate
\end{center}
\hrule
\bigskip

\begin{problem}
某消费者的效用函数为 $u(x,y)=x^{\alpha}y^{1-\alpha}$，其中
$0<\alpha<1$。价格 $p_x,p_y$ 和收入 $m$ 均为正。
求两种商品的需求量。
\end{problem}

\begin{solution}
效用对两种商品都严格递增，因此预算约束取等号：
$p_xx+p_yy=m$。内点最优时，一阶条件给出
\[
  \frac{\alpha}{1-\alpha}\frac{y}{x}=\frac{p_x}{p_y}.
\]
结合预算约束，可得
\[
  x^*=\frac{\alpha m}{p_x},
  \qquad
  y^*=\frac{(1-\alpha)m}{p_y}.
\]
\end{solution}

\begin{problem}
证明题目 1 得到的需求量恰好用完消费者的收入。
\end{problem}

\begin{solution}
\begin{proof}
将需求量代入总支出，可得
\[
  p_xx^*+p_yy^*
  =p_x\frac{\alpha m}{p_x}
   +p_y\frac{(1-\alpha)m}{p_y}
  =\alpha m+(1-\alpha)m=m.
\]
因此预算恰好用完。
\end{proof}
\end{solution}

\end{document}
